% !PRINTED 666
\weberpage{666}
%<<CONT
$z$ has a finite value. In order to determine the others as well, one must further make use of a second function, say $z_1 = 1 : z$, which has the finite value 0 at the still missing points, which are present in finite number.

\gesperrt{The totality of the points forms the absolute Riemann surface\textsuperscript{1)}.}

\begin{center}
\textbf{§ 188. Order Numbers.}
\end{center}

1. \gesperrt{If $\mathfrak{P}$ is a point, then to every function $\eta$ that receives at this point a finite value different from zero we assign the order 0 in $\mathfrak{P}$.}

2. \gesperrt{If $\omega$ runs through the totality of the functions vanishing in $\mathfrak{P}$, and $\varrho$ is one of these functions for which all}
\[
\frac{\omega}{\varrho}
\]
\gesperrt{have finite values, then $\varrho$ has the order 1.}

One also says that $\varrho$ becomes infinitely small in $\mathfrak{P}$ in the first order.

If $\dfrac{\varrho'}{\varrho}$ is neither zero nor infinity, then by this $\varrho'$ likewise has the order 1. Conversely, if $\varrho$ and $\varrho'$ both have the order number 1, then $\varrho : \varrho'$ and $\varrho' : \varrho$ are neither zero nor infinity.

3. \gesperrt{If $\varrho$ is of the first order and $\omega$ is any function in $\Omega$, then $\omega$ has the order number $n$ when $\varrho^{-n}\omega$ becomes neither zero nor infinite in $\mathfrak{P}$.}

This determination of the order $n$ is independent of which function of the first order we have taken. If $n$ is positive, $\omega$ becomes zero; if $n$ is negative, $\omega$ becomes infinite; and if $n = 0$, then $\omega$ has in $\mathfrak{P}$ a finite value different from zero.

\medskip\noindent\rule{4cm}{0.4pt}\par
\textsuperscript{1)} According to \gesperrt{Riemann's} theory of algebraic functions, every point of the closed many-sheeted surface corresponds to a point in our sense. If $z$ is the independent variable, then the \gesperrt{Riemann} surface is spread out in $n$ sheets over the $z$-plane.
As an absolute \gesperrt{Riemann} surface one may regard any surface to which the totality of those many-sheeted surfaces can be referred in a one-to-one and continuous manner.
